\question
«Умный дом» оснащён тремя датчиками: $A$ -- датчик открытия входной двери ($A=1$, если дверь открыта), $B$ -- датчик открытия окна ($B=1$, если окно открыто), $C$ -- датчик движения в комнате ($C=1$, если движение обнаружено). Сигнализация срабатывает ($F=1$), если дверь открыта, а окно закрыто; либо окно открыто и обнаружено движение.

\begin{enumerate}
\item Составьте таблицу истинности и найдите СДНФ функции $F(A,B,C)$, а также её вектор значений (используя прямую нумерацию: первая переменная -- старший разряд).
\item Найденный в пункте 1 восьмибитный вектор значений разбили на два блока по 4 бита (первый -- для первых четырёх строк, второй -- для последних четырёх) и передали по каналу связи, закодировав каждый блок кодом Хэмминга. При передаче в одном из блоков произошла однократная ошибка. Принята последовательность:
\[
1101001\ \ 1010001
\]
Найдите и исправьте ошибку, восстановите переданный 8-битный вектор значений функции и убедитесь, что он совпадает с вектором, найденным в пункте 1.
\end{enumerate}
---------------

Автор -- Михаил Овчинников